\chapter*{Lampiran}
\addcontentsline{toc}{chapter}{LAMPIRAN}
\markboth{Lampiran}{Lampiran}
\section*{Lampiran 1. Instrumen Angket Pra Penelitian Peserta Didik}
\markright{Lampiran 1. Instrumen Angket Pra Penelitian Peserta Didik}
\addcontentsline{toc}{section}{Lampiran 1. Instrumen Angket Pra Penelitian Peserta Didik}
\begin{center}
    \textbf{INSTRUMEN ANGKET PRA PENELITIAN PESERTA DIDIK}
\end{center}
\textbf{Petunjuk:}
\begin{enumerate}
    \item Isilah identitas kalian pada tempat yang telah disediakan.
    \item Perhatikan setiap pertanyaan dengan cermat dan seksama.
    \item Pilih jawaban yang sesuai dengan kepribadian anda.
    \item Seluruh pertanyaan wajib diisi.
    \item Hasil angket ini tidak mempengaruhi nilai dan dijaga kerahasiaannya.
\end{enumerate}
Identitas responden\\
Nama:\\
Kelas:\\
\begin{table}[H]
    \centering
    \begin{tabular}{|c|p{7cm}|c|p{3cm}|}
        \hline
        \textbf{No} & \textbf{Aspek} & \textbf{Bentuk Jawaban} & \textbf{Keterangan} \\
        \hline
        1. & Apakah pelajaran matematika sulit? & Pilihan opsi & 
        \begin{tabular}[t]{@{}l@{}}
            - Ya\\
            - Tidak
        \end{tabular} \\
        \hline
        2. & Apakah materi bangun ruang sisi datar sulit dipahami? & Pilihan opsi & 
        \begin{tabular}[t]{@{}l@{}}
            - Ya\\
            - Tidak
        \end{tabular} \\
        \hline
        3. & Saya lebih tertarik dengan pembelajaran praktik, daripada hanya dijelaskan oleh guru. & Pilihan opsi &
        \begin{tabular}[t]{@{}l@{}}
            - Ya\\
            - Tidak
        \end{tabular} \\
        \hline
        4. & Saya menyukai pendekatan pembelajaran yang membuat peserta didik menjadi lebih aktif. & Pilihan opsi &
        \begin{tabular}[t]{@{}l@{}}
            - Ya\\
            - Tidak
        \end{tabular} \\
        \hline
        5. & Saya menyukai pembelajaran dengan menggunakan teknologi & Pilihan opsi &
        \begin{tabular}[t]{@{}l@{}}
            - Ya\\
            - Tidak
        \end{tabular} \\
        \hline
    \end{tabular}
\end{table}

\newpage
\thispagestyle{empty}
\section*{Lampiran 2. Revisi Ahli Materi}
\markright{Lampiran 2. Revisi Ahli Materi}
\addcontentsline{toc}{section}{Lampiran 2. Revisi Ahli Materi}
\begin{center}
    {\large\bfseries SURAT VALIDASI}
\end{center}
\vspace{0.5cm}
\noindent Menerangkan bahwa yang bertanda tangan di bawah ini:
% \begin{table}[H]
%     \begin{tabular}{@{}ll@{}}
%         Nama    & : Wibowo Ramadhiyanto, S.Pd. \\
%         NBM     & : 1 411 355 \\
%     \end{tabular}
% \end{table}
\begin{table}[H]
    \begin{tabular}{@{}ll@{}}
        Nama    & : Rima Aksen Cahdriyana, M.Pd. \\
        NIDN     & : 0517049001 \\
    \end{tabular}
\end{table}
\noindent Telah melakukan validasi terhadap media pembelajaran, sebagai kelengkapan penelitian yang berjudul 
% \textbf{Pengembangan Media Pembelajaran Matematika Berbasis Simulasi 3D Menggunakan React dan Blender pada Materi Bangun Ruang Sis Datar}.
\textbf{PENGEMBANGAN MEDIA PEMBELAJARAN MATEMATIKA BERBASIS SIMULASI 3D MENGGUNAKAN REACT DAN BLENDER PADA MATERI BANGUN RUANG SISI DATAR}.

\noindent Yang disusun oleh:
\begin{table}[H]
    \begin{tabular}{@{}ll@{}}
        Nama    & : Toriq Afanudin \\
        NIM     & : 1900006105 \\
        Prodi   & : Pendidikan Matematika \\
        Fakultas & : Keguruan dan Ilmu Pendidikan \\
    \end{tabular}
\end{table}

\noindent Setelah memperhatikan media pembelajaran yang dikembangkan, ahli materi menyimpulkan bahwa media pembelajaran layak digunakan dengan revisi terlampir.

\noindent Dengan adanya validasi ini, diharapkan diperoleh kualitas media pembelajaran yang baik.\\[1.5em]
\begin{flushright}
\begin{minipage}{0.5\textwidth}
\centering
Yogyakarta, 15 Mei 2026\\[2.5cm]

Rima Aksen Cahdriyana, M.Pd.
\end{minipage}
\end{flushright}
\newpage
\thispagestyle{empty}
{\large\bfseries Lampiran: Revisi Ahli Materi}\\
\noindent Berikut adalah saran perbaikan dari ahli materi:
\begin{enumerate}[leftmargin=*]
    \item Kurang kata "satuan" dalam kolom input jawaban kubus satuan.
    \item Kalimat "kubus tersebut tidak terbentuk dari kubus satuan" kurang tepat secara konsep. Semua kubus sebenarnya bisa dibayangkan tersusun dari kubus satuan, meskipun tidak digambar secara eksplisit. Sebaiknya disampaikan "tidak ditampilkan dalam bentuk kubus satuan".
    \item Kata "representasi" kurang cocok untuk siswa SMP. Gantilah dengan "banyaknya kubus satuan yang menyusun suatu kubus menunjukkan volume kubus tersebut".
    \item Definisi prisma perlu diperbaiki karena berpotensi membingungkan siswa.
    \item Tambahkan satu paragraf sebelum definisi jaring-jaring. Misalnya: "Perhatikan kardus pada gambar di atas. Kardus tersebut berbentuk kubus. Untuk menghitung luas permukaan kubus, kita perlu mengetahui luas seluruh sisi-sisinya. Salah satu cara untuk melihat semua sisi kubus adalah dengan membukanya menjadi jaring-jaring."
    \item Tambahkan variasi jaring-jaring prisma lebih beragam.
    \item Menambahkan Tujuan Pembelajaran, Capaian Pembelajaran, dan Indikator Pencapaian Kompetensi pada bagian awal media pembelajaran.
\end{enumerate}

\newpage
\thispagestyle{empty}
\section*{Lampiran 3. Glosarium Istilah Teknologi Informasi (IT)}
\markright{Lampiran 3. Glosarium Istilah IT}
\addcontentsline{toc}{section}{Lampiran 3. Glosarium Istilah IT}
\begin{center}
    {\large\bfseries GLOSARIUM ISTILAH TEKNOLOGI INFORMASI}
\end{center}
\vspace{0.5cm}

\noindent Berikut adalah daftar penjelasan dari beberapa istilah sulit atau teknis, khususnya di bidang Teknologi Informasi (IT), yang digunakan dalam skripsi ini:

\begin{description}[leftmargin=1cm, style=nextline]
    \item[\textbf{Augmented Reality (AR)}] Teknologi yang menggabungkan benda maya dua dimensi maupun tiga dimensi ke dalam lingkungan nyata, lalu memproyeksikan benda-benda maya tersebut ke dalam waktu nyata (\textit{real-time}).
    \item[\textbf{Blender}] Perangkat lunak sumber terbuka (\textit{open-source}) yang digunakan untuk membuat dan merancang objek grafis 3D, meliputi pemodelan, animasi, hingga \textit{rendering}.
    \item[\textbf{Continuous Deployment}] Proses otomatis yang mempublikasikan dan menerapkan pembaruan kode aplikasi langsung ke peladen (server) secara langsung setelah diperbarui di repositori, tanpa perlu diunggah secara manual.
    \item[\textbf{Format .glb / .gltf}] Format berkas (\textit{file format}) 3D yang ringan dan efisien, dirancang khusus untuk memuat dan mentransmisikan model 3D beserta animasi dan teksturnya di \textit{platform web}.
    \item[\textbf{MyWebAR}] Sebuah \textit{platform} berbasis \textit{web} yang memungkinkan penggunanya untuk membuat dan membagikan pengalaman \textit{Augmented Reality} langsung melalui \textit{browser} tanpa mengharuskan pengguna menginstal aplikasi tambahan di ponsel.
    \item[\textbf{React}] Sebuah pustaka (\textit{library}) JavaScript sumber terbuka yang dikembangkan oleh Meta, digunakan secara khusus untuk membangun antarmuka pengguna (UI) situs \textit{web} agar lebih responsif.
    \item[\textbf{Single Page Application (SPA)}] Pendekatan pengembangan aplikasi \textit{web} di mana sebuah situs hanya memuat satu halaman HTML tunggal dan secara dinamis memperbarui kontennya saat pengguna berinteraksi, tanpa harus memuat ulang (\textit{reload}) halaman secara keseluruhan.
    \item[\textbf{Supabase}] Platform \textit{Backend-as-a-Service} (BaaS) sumber terbuka yang digunakan sebagai alternatif *Firebase*, dalam penelitian ini digunakan khusus sebagai \textit{storage} untuk menyimpan dan melayani berkas aset 3D.
    \item[\textbf{Three.js}] Pustaka (\textit{library}) JavaScript yang sangat populer digunakan untuk membuat dan menampilkan animasi grafis 3D pada \textit{browser web} dengan memanfaatkan teknologi WebGL.
    \item[\textbf{Vercel}] Layanan *cloud platform* yang digunakan untuk melakukan penyebaran (\textit{deployment}) atau \textit{hosting} situs \textit{web} secara instan langsung dari repositori *GitHub*.
    \item[\textbf{WebGL (Web Graphics Library)}] API (Antarmuka Pemrograman Aplikasi) JavaScript untuk merender grafis 2D maupun 3D interaktif yang kompatibel dengan \textit{browser web} modern tanpa memerlukan perangkat lunak tambahan (*plug-in*).
\end{description}
